\subsection{Residual capacity as a screening constraint}\label{sec:residual}

Deletion containment gives a convenient sufficient condition for feasibility,
but excludes profitable changes in ownership. We instead fix a complete
mechanism for bidder 1 and optimize bidder 2 against the residual allocation.
For a jointly measurable capacity function $r:[0,1]^4\to[0,1]^2$, let
\[
\mathcal V(r)=\sup\{\operatorname{Rev}_2(M_2):
 M_2\text{ is pointwise DSIC/IR and }x_2\le r\}.
\]
The mechanisms in this definition remain jointly measurable, with integrable
payments. Grouping feasible pairs by their first component gives exactly
\[
\OPT=\sup_{M_1}\{\operatorname{Rev}_1(M_1)+\mathcal V(1-x_1)\}.
\]
This identity does not assert existence of a best response or invoke a
measurable selection of pointwise optimizers.

\begin{proposition}[Residual value and supporting measures]\label{prop:residual-value}
The functional $\mathcal V$ is nondecreasing and concave. On a fixed opponent
fiber, suppose a finite nonnegative capacity-price measure $\pi$ satisfies
$R(u,a)\le\langle\pi,a\rangle$ for every admissible conditional mechanism,
and a feasible candidate at $r^*$ attains
$R^*=\langle\pi,r^*\rangle$. Then that candidate solves the full randomized
conditional problem and, for every residual for which the pairing is defined,
\[
V(r)\le V(r^*)+\langle\pi,r-r^*\rangle.
\]
\end{proposition}
\begin{proof}
Monotonicity is immediate. Mix any two admissible mechanisms using one
exogenous coin with constant probability $\theta$. All pointwise incentive,
IR and capacity inequalities are preserved, as are measurability and
integrability. Apply this to arbitrarily near-optimal competitors and pass
to the supremum. For the second assertion,
$R\le\langle\pi,a\rangle\le\langle\pi,r\rangle$, and equality at $r^*$
gives the claim. No own-report-dependent mixture is used.
\end{proof}

Concavity is the standard convex-mixture observation. The substantive step
below is an explicit support for residuals with binding report-line
constraints. These supports price information rents as well as item supply.
They are conditional certificates; a single measure supporting both bidders
simultaneously is a separate, unresolved requirement.

A report line has zero probability under the uniform distribution, but DSIC
links utility at that line to utility at all other reports. Its capacity
constraint can therefore restrict implementable allocations on a region of
positive probability. A line price values this incentive restriction; it
does not represent positive-probability sales on the line itself. All residuals
in the two conditional theorems below are Borel, so their singular pairings are
well-defined. The DSIC comparison class remains that of Section 3.

\paragraph{A sufficient common-certificate interface.}
On the full report space $\Omega=[0,1]^4$, suppose nonnegative finite measures
$\Pi_j$, one per physical item, have well-defined pairings with every
admissible allocation and satisfy, for each bidder's complete DSIC/IR mechanism,
\[
 R_i(M_i)\le\sum_{j=1}^2\int_\Omega x_{ij}\,d\Pi_j.
\]
If both inequalities are equalities at a joint candidate and
$\sum_j\int(1-x_{1j}^*-x_{2j}^*)\,d\Pi_j=0$, summing and using capacity gives
\[
 R_1(M_1)+R_2(M_2)\le\sum_j\Pi_j(\Omega)=R_1^*+R_2^*.
\]
This is a sufficient interface, not an existence, completeness or strong-duality
claim. Conditional supports must first be assembled on the same full report
space, with both revenue inequalities and their pairing domains justified.
Throughout the following historical construction, the final mechanism means
the V4.6 mechanism $M_J$. The active mechanism $M_*$ is specified later
in Section~\ref{sec:refined-mechanism}.
The strict joint gain below rules out a matching global certificate for its
reference before the release. A matching certificate for the final mechanism
is still unknown.

\subsection{A diagonal-hole capacity certificate}\label{sec:diagonal-support}

Use aligned own coordinates $(x,y)$, where item 1 is scarce. Put $A=2/3$,
$k=C-B<1/2$, $0\le k$, and
\[
 B=C-k\le A,\qquad C_0=5/6+3k^2/4,\qquad C\ge C_0.
\]
The candidate menu has prices $(0,A,B,C)$ and utility
$u^*=\max(0,x-A,y-B,x+y-C)$. Assume its complete allocation is feasible
against a Borel residual $r^*$ and that (the first two conditions are
needed only when $m=2(C-C_0)>0$)
\begin{equation}\label{eq:diagonal-traces}
\begin{split}
r^*_1(x,0)&=0\quad(0<x<1/2),\\
r^*_2(1/2,y)&=0\quad(0<y<C-1/2),\\
r^*_1(x,1)&=0\quad(0<x<k)\quad\text{if }B<A.
\end{split}
\end{equation}
Define $\ell(x)=B$ below $k$ and $\ell(x)=C-x$ on $(k,A)$;
let $D_0=\{x<A,y<\ell(x)\}$. Set $f=3\ell-2$ below $A$ and $f=1$
above $A$. Direct integration gives
\[
m=\int_0^1 f=2(C-C_0)\ge0,\qquad 3|D_0|=m+1.
\]
The function $F_m(x)=\int_x^1f(s)\,ds-m\mathbf1_{x<1/2}$ is nonnegative.
Indeed, below the anchor it equals $-\int_0^x f$, while above the anchor
its tail is nonnegative because $f\le0$ below $A$ and $f=1$ above $A$.

\begin{theorem}[Anchored full screening support]\label{thm:diagonal}
Under these hypotheses the candidate solves the complete randomized
conditional problem. A supporting capacity measure is
\[
\begin{aligned}
\pi_1={}&3(x-A)_+\,dx\,dy+F_m(x)\,dx\,\delta_1(dy)
                +m\mathbf1_{0<x<1/2}\,dx\,\delta_0(dy),\\
\pi_2={}&3(y-\ell(x))_+\mathbf1_{x<A}\,dx\,dy
                +m\delta_{1/2}(dx)\,dy.
\end{aligned}
\]
For every Borel residual $r$ and randomized DSIC/IR competitor with selected
allocation $a\le r$,
\begin{equation}\label{eq:diagonal-gap}
\langle\pi,r\rangle-R
 =\langle\pi,r-a\rangle+3\int_{D_0}u-mu(0,0)\ge0.
\end{equation}
\end{theorem}
\begin{proof}
Uniform revenue is
$R=\int u(1,y)\,dy+\int u(x,1)\,dx-3\int u$.
Horizontal envelopes on $x>A$ and vertical envelopes above $\ell(x)$
give the displayed volume prices and the residual top pairing
$\int f(x)u(x,1)\,dx$. Anchoring it at $1/2$ gives
\[
\int f(x)u(x,1)\,dx
 =mu(1/2,1)+\int F_m(x)a_1(x,1)\,dx.
\]
Expand $u(1/2,1)$ along the bottom edge and then the vertical line at $1/2$.
This produces the remaining two prices and \eqref{eq:diagonal-gap}.
Monotonicity and IR imply $u\ge u(0,0)\ge0$, so the non-capacity slack is
at least $u(0,0)$ by $3|D_0|=m+1$.

At the candidate this slack vanishes. Every positive volume price charges a
unit marginal. The top price below $k$ is zero if $B=A$ and is saturated by
the last condition in \eqref{eq:diagonal-traces} otherwise. The other two
conditions saturate the bottom and vertical charges. All line endpoints
have zero price mass, but retain their complete allocation rules.
\end{proof}

This proof uses the actual tangential allocation on each priced line.
Pointwise DSIC makes that component monotone and equal to the convex trace's
derivative almost everywhere on the line. It therefore covers arbitrary
selected subgradients and allocations with an infinite range. It neither
normalizes $u(0,0)$ to zero nor restricts competitors to a finite menu.
At $B=A$, $C=b_0=(4-\sqrt2)/3$, the mass $m$ is zero and no occupied
trace is required. This recovers the full-capacity single-buyer support.

\subsection{A lottery-junction capacity certificate}\label{sec:lottery-support}

Set $c=137/500$, $q=113/500$, $T=613/750$, and $A<t<T$. Write
\[
\begin{gathered}
\alpha=3t-2,\quad
\delta=q+3q^2/4-3qt/2+\alpha^2/16,\quad
\beta=\frac{\alpha}{\alpha+2\delta},\\
B=1/2+\delta,\quad C=t+c+\delta,\quad k=t-q,\quad
\ell_*=\delta+\alpha/2,\quad Y=c+\ell_*,\quad j=1-t/2.
\end{gathered}
\]
Here $0<k<j<A<t<1$, $c<Y<B<A$ and $0<\beta<1$. Offer, in order,
\[
(0,0;0),\ (0,1;B),\ (1,0;t),\ (1,1;C),\
(1,\beta;t+\beta c),
\]
choosing the first utility maximizer. Assume feasibility at $r^*$ and
\begin{equation}\label{eq:lottery-traces}
r^*_1(x,1)=0\ (x<k),\qquad
r^*_1(x,c)=0\ (A<x<t).
\end{equation}
Define $z(y)$ by $t$ below $c$, $t-\beta(y-c)$ on $(c,Y)$, and $A$
above $Y$. Define $e(x)$ by $B$ below $k$, $C-x$ on $(k,j)$, and $Y$
on $(j,A)$. Horizontal integration on $G=\{x>z(y)\}$ and vertical
integration on $H=\{x<A,y>e(x)\}$ partition the positive-utility region
up to volume-zero faces. Put $D_0=[0,1]^2\setminus(G\cup H)$; the
candidate has zero utility \emph{almost everywhere} on $D_0$. This statement
about a volume integral does not discard pointwise capacity or line prices.

Let $f=3e-2$ below $A$ and $f=1$ above $A$. The displayed formula for
$\delta$ implies $\int f=0$, and $F(x)=\int_x^1 f\ge0$. Set
$\lambda=\alpha(c+\ell_*/4)$ and let $w(x)$ equal zero below $A$,
$3t(c+\ell_*/4)$ on $(A,t)$, and $\lambda$ above $t$.

\begin{theorem}[Full randomized lottery support]\label{thm:lottery}
Under these hypotheses the five-option candidate is optimal over all
randomized conditional mechanisms. A supporting measure is
\[
\begin{aligned}
\pi_1&=3(x-z(y))_+\,dx\,dy+F(x)\,dx\,\delta_1(dy)
                         +w(x)\,dx\,\delta_c(dy),\\
\pi_2&=3(y-e(x))_+\mathbf1_{x<A}\,dx\,dy.
\end{aligned}
\]
For every Borel residual $r$ and competitor with $a\le r$,
\begin{equation}\label{eq:lottery-gap}
\langle\pi,r\rangle-R
 =\langle\pi,r-a\rangle+T_g+\lambda M_a+S_u\ge0.
\end{equation}
\end{theorem}
\begin{proof}
Let $g(y)=u(1,y)$ and let $\nu$ be its nonnegative second-derivative
measure. The horizontal and vertical envelopes leave top coefficient $f$
and right coefficient $\sigma(y)$, equal to $\alpha$ below $c$,
$\alpha-3\beta(y-c)$ on $(c,Y)$ and zero above $Y$.
The top coefficient integrates to the stated top-line price. Convexity on
the whole lottery interval gives
\[
\begin{split}
T_g&=\lambda g(c)-\int\sigma(y)g(y)\,dy\\
&=\alpha\int_0^c[g(c)-g(y)]\,dy
 +\frac\beta2\int_0^{\ell_*}s(\ell_*-s)^2\,d\nu(c+s)\ge0.
\end{split}
\]
This follows from the exact hinge coefficient
$\int_s^{\ell_*}(\alpha-3\beta z)(z-s)\,dz
=-\beta s(\ell_*-s)^2/2$ and $\beta\ell_*=\alpha/2$.
Next, monotonicity of $a_1(\cdot,c)$ gives
\[
M_a=\frac{t}{t-A}\int_A^t a_1(x,c)\,dx-
                     \int_0^t a_1(x,c)\,dx\ge0.
\]
Expand $g(c)$ along this line to obtain its price $w$.
Finally,
$S_u=3\int_{D_0}u-\lambda u(0,c)\ge0$ because
$D_0$ contains $[0,k]\times[c,B]$ up to volume-zero faces and
\[
3k(B-c)-\lambda\ge q(1-3q/2)>0.
\]
These identities yield \eqref{eq:lottery-gap} without removing origin
utility. At the candidate the trace is constant below $c$, affine on
$(c,Y)$, and all three non-capacity slacks vanish. Feasibility and
\eqref{eq:lottery-traces} saturate the prices. The safe-item price is zero
on the genuine lottery cell, consistent with its fractional allocation.
\end{proof}

The full slack identity rules out higher conditional revenue from additional
lotteries or more general randomized allocations, including scarce marginals
below one. Its hinge term controls nonlinear convex trace changes over the
whole lottery interval. The finite menu is an equality case of an unrestricted
bound; neither uniqueness nor a classification of all equality cases is asserted.

% Author revision: insert as an appendix or a subsection of the construction.
% Requires amsmath, amssymb, booktabs, array and hyperref (already in manuscript).
\subsection{Specification of the retained V4.6 mechanism and its revenue}
\label{sec:active-reference-specification}

The retained V4.6 lower-bound mechanism is $M_J$, including the final joint price
release. Its reference before the new modifications is $M_{4.5}$. This
specification fixes their parameters, order of composition and exceptional
reports. All source paths below are relative to
\path{certificate/joint_residual_screening_lower_bound/source/} in the
archive. A profile is ordered as
$\omega=(w_1,w_2,v_1,v_2)\in[0,1]^4$: $w$ is bidder 1's report and $v$ is
bidder 2's report, each in physical item order. A menu option $(x;p)$ gives
expected allocation $x\in[0,1]^2$, payment $p$, and utility $w\cdot x-p$
or $v\cdot x-p$ for the respective bidder.
Masks $0,1,2,3$ mean empty, item 1, item 2 and both items.

\paragraph{Constants and retained tariffs.}
Put
\[
\begin{gathered}
 a=\frac{159}{250},\quad b=\frac{91}{100},\quad c=b-a=\frac{137}{500},\\
 s_0=\frac{1137}{1000},\quad s=\frac{142}{125},\quad
 d=s-a=\frac12,\quad q=s-b=\frac{113}{500}.
\end{gathered}
\]
The earlier deterministic construction uses $s_0$, hence
$d_0=501/1000$ and $q_0=227/1000$. This historical mechanism changes the shared
allocation cost to $s=s_0-1/1000$ and retains the V3 price increments as
functions of the opponent report. In particular, the constants in those
frozen functions are not replaced by $d,q$.
For an opponent report $u$, define
\[
\begin{split}
 H(u)&=\max\{0,u_1-a,u_2-a,u_1+u_2-b\},\\
 m^r(u)&=(0,H(u)+\min\{a,r-u_2\},
                 H(u)+\min\{a,r-u_1\},H(u)+b).
\end{split}
\]
Let $P^0(u)$ be the four prices returned by
\path{V3/verifier/joined_threshold.py:menu(u)}, using
\path{V3/verifier/baseline_mechanism.py:menu(u)['prices']} when the first
function returns \texttt{None}. These are explicit piecewise algebraic
functions: the joined rule applies on
$d_0\le t\le1$, $0\le\rho\le b-\min(t,a)$, where
$t=\max(u)$, $\rho=\min(u)$; its successive branches use the tests
$t\le a$, $(t-q_0)^2\le r_0$, $(t-q_0)^2<r_1$, and $t\le2/3$, with
$r_0=b^2-(2/3+c^2)$ and $r_1=4c/3-2/9$ exactly as in that function.
The remaining affine branch is
\path{V3/verifier/stationary_s.py:prices_for_t(t)}.
The finite tariff data are
\path{V3/certificate/baseline_mechanism.json}: 20 common rows, 41 bundle
rows and 8 item rows, in their stored order. The literal first-match and
endpoint rules are \path{baseline_mechanism.py:fee_rows} and
\path{baseline_mechanism.py:menu} in the same verifier directory.
Thus the retained prices used by both bidders are exactly
\begin{equation}\label{eq:active-retained-prices}
 P^s(u)=m^s(u)+P^0(u)-m^{s_0}(u).
\end{equation}
This is \path{V4_02/verifier/split_cost_candidate.py:retained_menu}
at \texttt{CANDIDATE\_THETA}$=-1/1000$.

\paragraph{Shared selection and the reference $M_{4.5}$.}
First maximize reported total value minus cost over the ordered mask pairs
\[
 (0,0),(1,0),(2,0),(3,0),(0,1),(0,2),(0,3),(1,2),(2,1),
\]
with respective costs $0,a,a,b,a,a,b,s,s$, choosing the first maximizer.
This is the retained source's literal \texttt{base\_outcome\_order}, including
$(3,0)$ before $(0,1)$; it fixes the tie convention for this construction.
For each bidder maximize utility at $P^s(u)$: choose empty if the maximum
is zero; otherwise retain that bidder's shared mask if it maximizes utility;
otherwise choose the smallest maximizing mask contained in the shared mask.
This selection rule is part of the mechanism.

Next override bidder 2's menu when
$w\in Q=\{\max(w)\le A,\ w_1+w_2\le b\}$, where $A=2/3$ and
$b_0=(4-\sqrt2)/3$. If $t=\max(w)\le d$, its singleton and bundle prices
are $(A,A,b_0)$, with mask priority $(0,1,2,3)$. If $t>d$, put
$k=t-q$, $C_0=5/6+3k^2/4$ and use aligned prices
$(A,C_0-k,C_0)$ for the scarce singleton, safe singleton and bundle,
with priority (empty, safe, scarce, bundle). The scarce item is the item
at the opponent's larger coordinate. The source predicate is
\texttt{u[0] > u[1]} for scarce item 1 and item 2 otherwise; equal
coordinates cannot occur in a constrained $Q$ row because $t>d$ and
$\sum u\le b<2d$. These closed-$Q$ rules are
\path{split_cost_candidate.py:screening_menu}; the resulting evaluator is
\path{V4_02/verifier/split_cost_candidate.py:candidate}.

Finally, for each bidder whose opponent satisfies
$A<t<T=613/750$ and $u_1+u_2\le b$, replace its menu by
\begin{equation}\label{eq:active-lottery-menu}
 (0,0;0),\ (0,1;d+\delta),\ (1,0;t),\ (1,1;t+c+\delta),\
 (1,\beta;t+\beta c),
\end{equation}
in the displayed priority, where
\[
 \delta=q+\frac{3q^2}{4}-\frac{3qt}{2}+\frac{(3t-2)^2}{16},
 \qquad \beta=\frac{3t-2}{3t-2+2\delta}.
\]
The first aligned coordinate is the scarce item; reverse allocation
coordinates when $u_1<u_2$. This gives $M_{4.5}$, evaluated by
\path{V4_5/verifier/residual_lottery.py:mechanism}.
The endpoints $t=A,T$ retain the preceding rule; the face $\sum u=b$
is included in the lottery rule.

\paragraph{Ordered modifications defining $M_J$.}
Apply the following operations to $M_{4.5}$, retaining every earlier row
unless it is explicitly replaced. First, for both bidders extend
\eqref{eq:active-lottery-menu} to $A<t<T$, $b-t<\rho<c$; retain the
preceding rule on $\rho=c$. Second, for bidder 1's opponents in
$F=\{\max(v)\le d,\ \sum v\le b_0\}$ use prices $(A,A,b_0)$.
For opponents in $G=\{\max(v)\le d,\ b_0<\sum v<b\}$ use
$(a,a,\sum v)$. Both use mask priority $(0,1,2,3)$.
Third, on bidder 2's $Q$ rows replace the bundle price by
$\max(b_0,\sum w)$ if $\max(w)\le d$.
On a constrained row, if $\sum w>C_0$, replace the aligned prices by
$(A,\sum w-k,\sum w)$; otherwise retain its preceding $Q$ menu.
Keep the $Q$ priorities above. The boundaries $\sum v=b_0,b$ retain their
preceding bidder-1 rules, and equality $\sum w=C_0$ retains the earlier
constrained row. These operations, including the intermediate bundle-only
response before the coupled response, are the successive
\texttt{mechanism} functions in
\path{V4_6/verifier/outer_lottery_strip.py},
\path{free_bidder_one.py}, \path{outer_bundle_exchange.py} and
\path{outer_bundle_reoptimized.py}, all in that verifier directory.
Call the resulting mechanism $M_6$.

Set $\varepsilon=9/10000$ and use the closed rectangles
\[
 W=[7/10,71/100]\times[0,1/5],\qquad
 S=[7/10-4\varepsilon,71/100]\times[c-2\varepsilon,c+4\varepsilon].
\]
When $v\in S$, charge bidder 1 an additional $\varepsilon$ on every
nonempty option of its $M_6$ menu: choose empty when its previous maximum
utility is at most $\varepsilon$, and otherwise retain its selected option.
When $w\in W$, offer bidder 2 the five options in
\eqref{eq:active-lottery-menu}, lowering only the lottery payment by
$\varepsilon$, with the same first-maximizer priority. This last operation
uses the displayed physical orientation only. The final evaluator is
\path{V4_6/verifier/price_joint_reallocation.py:mechanism(profile)}.
Its \texttt{allocations} and \texttt{payments} fields define $M_J$.
The real mechanism interprets every displayed and source-defined branch
over the reals; the executable evaluates rational reports using exact
rational and quadratic comparisons. No floating-point extension or numerical
tie tolerance defines the mechanism.

\paragraph{Exact revenue.}
For each specified evaluator $E$, let $p_i^E(w,v)$ be its selected payment.
With four independent uniform values, define the exact number
\begin{equation}\label{eq:active-payment-integral}
 \mathcal R(E)=\int_0^1\!\int_0^1\!\int_0^1\!\int_0^1
 [p_1^E(w_1,w_2,v_1,v_2)+p_2^E(w_1,w_2,v_1,v_2)]
 \,dv_2\,dv_1\,dw_2\,dw_1.
\end{equation}
In particular $R_{4.5}=\mathcal R(M_{4.5})$ and
$R_J=\mathcal R(M_J)$. The following table fixes the reduced-integral
entrypoints; it does not define revenue through a saved decimal.

\begingroup\small
\renewcommand{\arraystretch}{1.18}
\begin{center}
\begin{tabular}{@{}>{\raggedright\arraybackslash}p{0.18\linewidth}>{\raggedright\arraybackslash}p{0.76\linewidth}@{}}
\toprule
Quantity & Exact definition and reduction in the source directory\\
\midrule
Historical base $R_{02}$ &
$\mathcal R(\texttt{split\_cost\_candidate.candidate})$.
\path{V4_02/verifier/split_cost_revenue.py:calculate} reduces its change
from V3.1 using \texttt{BASE\_D}, \texttt{base\_Q}, \texttt{inner\_Q}
and \texttt{moments}; V3.1 is
\path{V3_1/verifier/constrained_candidate.py:mechanism}.\\
$R_{4.5}-R_{02}$ &
$4\int_A^{t_o}(b-t)g_1(t)\,dt+
 4\int_{t_o}^{T}(b-t)g_2(t)\,dt$;
\path{V4_5/verifier/independent_lottery.py:reconstruct}.\\
$R_6-R_{4.5}$ &
$\int_{[0,1]^4}\sum_i(p_i^{M_6}-p_i^{M_{4.5}})$;
\path{V4_6/verifier/independent_revenue.py:calculate}
returns \texttt{refined\_addition\_coefficients} in the basis
$(1,\sqrt2,\sqrt{23})$.\\
$J=R_J-R_6$ &
$\int_{[0,1]^4}\sum_i(p_i^{M_J}-p_i^{M_6})$;
\path{V4_6/verifier/price_joint_revenue.py:calculate}
reduces this integral to a rational constant and three logarithms.\\
\bottomrule
\end{tabular}
\end{center}
\endgroup
Here $\delta_o(t)=1/2-c+3q_0^2/4-3q_0t/2$,
$t_o=(1/2-c+3q_0^2/4)/(3q_0/2)$,
$g_1=(\delta-\delta_o)^2+\delta_o(3t-2)^2/8$ and $g_2=\delta^2$.
The V4.02 reduction uses its explicitly integrated polynomial on
$[s_0,s]$ and the fixed tariff moments; its historical base ultimately
shares the earlier affine-polytope revenue calculation.
\path{V4_6_archive_audit/revenue/fresh_exact_gap.py:algebraic_increment}
and \texttt{joint\_increment} independently reconstruct the two new
V4.6 additions. They retain the V4.5 base enclosure and therefore do not
constitute a second integration of the entire historical base.
The rational radical and logarithm bounds in that file and
\path{V4_6/verifier/consolidated_bounds.py:verify} give
\[
 0.8758198541484224553460<R_J<0.8758198541484224553461.
\]
This definition and enclosure concern the constructed mechanism. They do
not assert unrestricted optimality or a matching full-auction certificate.


\subsection{Complete joint reallocations}\label{sec:joint-mechanism}

Throughout this subsection let $\mathbf w=(w_1,w_2)$ be bidder 1's report
and $\mathbf v=(v_1,v_2)$ bidder 2's, in physical item order. In the formulas
below $w$ and $v$ abbreviate these vectors, replacing the scalar-coordinate
notation of Section 3. A set of opponent reports indexes a complete own-type
menu. Lottery rows put the opponent's high item first and are rotated back to
physical items after the menu is defined; their equal-coordinate convention
uses item 1 first. The constrained Q rule retains its specified scarce-item
selection. The final rectangles W and S use physical item 1 as the high item and
are not rotated. Thus $((.45,.45),(.44,.44))$ means $(\mathbf w,\mathbf v)$,
whereas a final lottery change is indexed by $\mathbf w\in W$ and may require
a fee in the opposing menu indexed by $\mathbf v\in S$.

The reference $M_{4.5}$ is specified in
Section~\ref{sec:active-reference-specification}, with its finite menu rules,
rational tables and algebraic tests in the source-bound
certificate \path{joint_residual_screening_lower_bound}. Its complete
predecessor proofs and exact revenue integrals are included in that package.
The following modifications are applied in order. All new menu selections
have the stated fixed priority; unmodified rows retain their complete
predecessor maximizer selection, which need not be independent of the own
report. Selecting any utility-maximizing option from a fixed opponent menu
is sufficient for pointwise DSIC.

First extend both lottery menus to
$E^+=\{A<\max(v)<T,\min(v)<c\}$. Retain the predecessor on the
$\min(v)=c$ face. The retained competing bidder cannot take its low item
there, and the two occupied traces in \eqref{eq:lottery-traces} remain in
place. Theorem~\ref{thm:lottery} therefore solves both corresponding inner
problems at this reference residual.

Second, bidder 2 receives nothing at every opponent report whenever its
own type lies in the closed set
\[
F=\{\max(v)\le1/2,\ v_1+v_2\le b_0\},\qquad
b_0=(4-\sqrt2)/3.
\]
On opponents in F, replace bidder 1's constant menu $(a,a,b)$ by the full
single-buyer optimal menu $(A,A,b_0)$. Here $a=159/250$, $b=91/100$.
Capacity is entirely free on these fibers; zero-utility ties select empty.

Third, on $G=\{\max(v)\le1/2,b_0<v_1+v_2<b\}$, bidder 1 offers bundle
price $v_1+v_2$, retaining singleton prices $(a,a)$. Bidder 2's complete
response changes on $Q=\{\max(w)\le A,\sum w\le b\}$. For $\max(w)\le1/2$
its prices are $(A,A,\max(b_0,\sum w))$. In a constrained row write
$t=\max(w)>1/2$, $k=t-q$, $C_0=5/6+3k^2/4$, and $z=\sum w$.
If $z>C_0$, the aligned prices are $(A,z-k,z)$; otherwise retain its earlier
Q menu. Moving the safe and bundle prices together preserves $C-B=k$.
Relative to the earlier feasible Q menu this only discards allocations.
A new bidder-1 bundle requires $\sum w>\sum v$, and the opposing bundle
threshold prevents double allocation. At $(w,v)=((.45,.45),(.44,.44))$,
ownership transfers strictly from bidder 2 to bidder 1, who pays $.88$.

On the reopened Q rows, F/G bundle purchases supply the bottom and vertical
conditions in \eqref{eq:diagonal-traces}; the top hole is inherited from the
previous capacity rule. Theorem~\ref{thm:diagonal} recertifies the whole
response. This includes all Q after combining unchanged and changed rows.
It also shows why increasing only the opposing bundle price was insufficient.
The complete constrained conditional revenue change is exactly
$-(z-C_0)^2$ at the coupled response.

Finally, the conditional lottery certificate identifies a consumed
junction that admits a profitable joint capacity release. Use one physical
orientation and set
\[
\varepsilon=9/10000,\quad
W=[.7,.71]\times[0,1/5],\quad
S=[.7-4\varepsilon,.71]\times[c-2\varepsilon,c+4\varepsilon].
\]
The rectangles are closed. On bidder-2 menus indexed by $w\in W$, lower
only the lottery payment by $\varepsilon$. On bidder-1 menus indexed by
$v\in S$, add $\varepsilon$ to every nonempty option. Choose empty whenever
the old maximal utility is at most $\varepsilon$; otherwise retain its old
maximizer and charge the additional fee. All other rows are unchanged.

\begin{proposition}[Pointwise joint release]\label{prop:joint-release}
The composed mechanism is pointwise DSIC, ex-post IR and jointly feasible.
Its allocations are jointly measurable and its payments are integrable. The final release has strictly positive
exact revenue increment $J$ given below.
\end{proposition}
\begin{proof}
Each change specifies a complete opponent-dependent menu before maximizing;
the empty option and full predecessor rules handle every exceptional report.
For the final change only a newly chosen lottery can introduce a capacity
conflict, since the fee contracts bidder 1's allocation. Writing $v=(x,y)$,
its comparisons with singleton, bundle and empty force
\[
c-\varepsilon/\beta\le y\le Y+\varepsilon/(1-\beta),\qquad
x+\beta(y-c)-t+\varepsilon>0.
\]
Throughout W, $3/5<\beta<3/4$ and
$x\ge1-t/2-4\varepsilon>1/2$. The F/G modifications are therefore absent
at every newly chosen lottery report. Bidder 1 can only take the scarce
singleton, at price $P\ge\max(x,x+y-c)$. If $t>P$, the displayed lottery
inequalities force $(x,y)\in S$ and
\[
t-P\le t-\max(x,x+y-c)<\varepsilon.
\]
The entry fee selects empty, releasing that item. Its low item is already
free because its value is at most $1/5<c$. Equalities use the explicit empty
priority. The earlier splices have the full same-orientation, opposite-
orientation and reverse-Q capacity checks in the certificate proof notes;
these checks include the excluded $\rho=c$ and $\sum v=b$ faces.
Finite menus with Borel parameters yield joint measurability and bounded
payments. Feasible marginals can be realized itemwise.
\end{proof}

This release preserves the full F and Q certificates. It changes some
Eplus menus and their binding residuals. A sufficient preserved Eplus scope
excludes S for bidder 1 and excludes
$\max(w)\in[.7-4\varepsilon,.71+\varepsilon]$ for bidder 2. The wider
second exclusion includes releases outside W. No conditional optimality is
asserted on excluded fibers, and no common auction support is obtained.

\begin{table}[htbp]
\centering\small
\begin{tabular}{@{}>{\raggedright\arraybackslash}p{.11\linewidth}>{\raggedright\arraybackslash}p{.09\linewidth}>{\raggedright\arraybackslash}p{.31\linewidth}>{\raggedright\arraybackslash}p{.35\linewidth}@{}}
\toprule
Opponent set & Bidder & Complete menu change & Full conditional support after release\\
\midrule
F & Both & Single-buyer menu $(A,A,b_0)$ & Yes, at full residual capacity\\
Q & 2 & Retained or coupled diagonal response & Yes, Theorem~\ref{thm:diagonal} and retained rows\\
G & 1 & Bundle price equal to opponent total & No new inner-optimality claim\\
$E^+$ & Both & Five-option lottery & Theorem~\ref{thm:lottery} only outside the exclusions above\\
W & 2 & Lottery payment minus $\varepsilon$ & No matching inner claim on the changed rows\\
S & 1 & Fee $\varepsilon$ on every nonempty option & No matching inner claim on the changed rows\\
\bottomrule
\end{tabular}
\caption{Report roles and final certificate scope. Sets index opponent menus;
F is contained in Q, and W is contained in one orientation of $E^+$. Entries
state certified scope, not a claim that the other fibers are suboptimal.}
\label{tab:conditional-scope}
\end{table}

\subsection{Exact new revenue and a sub-0.01 bracket}\label{sec:joint-revenue}

Let $R_{4.5}=\mathcal R(M_{4.5})$ be the exact payment integral
\eqref{eq:active-payment-integral}. Its displayed reduction equals the V4.02
reference integral plus
$472617707798758460108653/101062797120000000000000000000$.
The baseline is an exact mathematical integral, not its printed decimal.
Exact continuous-cell integration of all subsequent complete menus gives
\begin{equation}\label{eq:joint-revenue}
R_J=R_{4.5}+a_0+a_2\sqrt2+a_{23}\sqrt{23}+J,
\end{equation}
where
\[
\begin{gathered}
a_0=-\frac{1609541379668231661249547038251}
              {159173905464000000000000000000000},\qquad
 a_2=\frac{57567793}{7593750000},\\
a_{23}=-\frac{12167}{1328906250},
\end{gathered}
\]
and
\[
\begin{split}
J={}&\frac{6578986777818169914603}{10652675687500000000000000000000}\\
&+\frac{6248637}{282500000000000}\log\frac{161}{176}
-\frac{61575471}{282500000000000}\log\frac{613}{628}\\
&-\frac{9308601}{5000000000000000000}\log\frac{13}{10}.
\end{split}
\]
Here $J=0.00000000391612794214225366689702614\ldots$ and an independently
proved rational lower bound is $J\ge26902077489/12500000000000000000>0$.

All induced information-rent changes are included. In a lottery row,
\[
\Delta R_2=\frac{\alpha\ell_*}{4}\varepsilon
-\frac{\alpha}{4\beta(1-\beta)}\varepsilon^2
-\frac{1}{2\beta(1-\beta)^2}\varepsilon^3.
\]
The coefficient $\lambda=\alpha(c+\ell_*/4)$ in the certificate's bound
on $g(c)$ is not this row's revenue derivative. On the fixed lottery interval,
the price cut changes the convex-trace slack by exactly
\[
 \Delta T_g=\alpha\left(c\varepsilon-\frac{\varepsilon^2}{2\beta}\right),
 \qquad
 \left.\frac{d\Delta R_2}{d\varepsilon}\right|_{0+}
       =\lambda-\alpha c=\frac{\alpha\ell_*}{4}.
\]
Indeed, below $c$ the right-trace increase is
$(\varepsilon+\beta(y-c))_+$, whose integral is
$\varepsilon^2/(2\beta)$; on $(c,Y)$ it is constant, so the hinge-curvature
term is unchanged. Thus even a saturated conditional support must be combined
with its changing incentive slacks before predicting the revenue from a
complete menu deformation.

For a complete entry-fee row with no-sale area $D$ and bundle price $C$,
\[
\Delta R_1=(1-3D)\varepsilon-\frac32 C\varepsilon^2-\frac12\varepsilon^3.
\]
This does not use the earlier four-option lemma's assumption $C\le1$.
For the actual rows and $0\le h\le\varepsilon$, the internal menu junctions
keep their order, the no-sale area is $D(h)=D+Ch+h^2/2$, and the top and
right utility traces stay positive after the fee. The uniform revenue
identity then gives $R'(h)=-2+3(1-D(h))=1-3D(h)$; integration proves the
formula. For example, the lottery row $t=7/10$ has
$\delta=1727/62500$ and $C=31301/31250>1$, while $t+\varepsilon<1$ and
$B+\varepsilon<1$ keep the boundary traces positive. Its no-sale polygon
satisfies the same area identity.

For the selected $\varepsilon=9/10000$, integrating both complete menu
changes gives
\[
\begin{gathered}
 J_1=-\frac{24631898853429}{125000000000000000000000}<0,\\
 J_2=J-J_1=0.0000000041131831329696856\ldots>0.
\end{gathered}
\]
The gain therefore pays the actual negative revenue change of bidder 1.
A direct strict certificate, using $|W|=1/500$ and the complete fee-strip
bounds, is
\[
\begin{split}
J\ge{}&\frac1{500}\left(\frac{1213}{625000}\varepsilon
 -\frac{13}{60}\varepsilon^2-\frac{40}{3}\varepsilon^3\right)\\
&-(1/100+4\varepsilon)\left(\frac2{25}\varepsilon^2
 +\frac{903}{100}\varepsilon^3+3\varepsilon^4\right)\\
={}&\frac{26902077489}{12500000000000000000}>0.
\end{split}
\]
The whole lottery row contributes a linear term, whereas the compensating
fee is integrated over a strip whose width shrinks with the perturbation.
These formulas explain the certified fixed choice. We do not infer a
uniformly admissible perturbation family, an optimal step size, or an
unrestricted stationarity condition from this comparison.

The literal retained tariff chambers, including the distinct $y=c$ face,
are checked independently by rational demand polygons. Integration over
opponent reports gives the expression above. Radical bounds use integer
square-root enclosures; logarithm bounds use signed rational series
remainders. An independent reconstruction uses direct bivariate menu
integration and local pole expansions instead of the primary partial-
fraction linear system.

The resulting rational enclosure is
\[
0.8758198541484224553460<R_J<0.8758198541484224553461.
\]
Consequently
\[
0<\Ustream-R_J<0.0099384483702321295386<\frac1{100}.
\]
This establishes the lower claim in Theorem~\ref{thm:bracket}. It improves
the deterministic predecessor by $0.001214867745586187\ldots$ and the
V4.5 reference by $0.000565326708609134\ldots$. These are internal certified
improvements, not comparisons with the reported GemNet revenue.
